\documentclass[10pt,a4paper]{amsart}
\usepackage[margin=19mm]{geometry}
\usepackage{amsmath,amssymb,mathtools}
\usepackage[scr=rsfs,cal=euler]{mathalfa}
\usepackage{xfrac}
\usepackage[unicode,hidelinks,bookmarks=false]{hyperref}
\newcommand{\ie}{i.e.\ }
\newcommand{\cf}{cf.\ }
\newcommand{\resp}{respectively\ }
\newcommand{\Subsection}{\subsection}
\providecommand{\nobreakdash}{\nobreak\hbox{-}}
\setlength{\emergencystretch}{2em}
\multlinegap0pt
\newcommand*\samethanks[1][\value{footnote}]{\footnotemark[#1]}
\usepackage{tikz}
\usepackage{xcolor}
\usepackage[shortlabels]{enumitem}
\newtheorem{lemma}{Lemma}
\newcommand{\R}{\mathbb{R}} % real numbers
\newcommand{\de}{\mathrm{d}}
\newcommand{\jfunc}{\mathsf{J}}
\newcommand{\mc}[1]{\mathcal{#1}}
\newcommand{\norm}[1]{\left\|{#1}\right\|} % norm
\title{Is Memorization Helpful or Harmful? \\ Prior Information Sets the Threshold}
\author{Chen Cheng \and Rina Foygel Barber\samethanks}
\date{Source: arXiv:2602.09405v2 (CC BY 4.0)}
\begin{document}
\maketitle

\noindent\emph{Source and licence.} Is Memorization Helpful or Harmful? \\ Prior Information Sets the Threshold. Version arXiv:2602.09405v2 (CC BY 4.0), distributed by arXiv under the Creative Commons Attribution 4.0 International licence, http://creativecommons.org/licenses/by/4.0/, as recorded in the arXiv licence metadata for this exact version and shown on the abstract page https://arxiv.org/abs/2602.09405v2. Full licence notice: \texttt{licenses/arxiv-2602.09405-CC-BY-4.0.txt}.\par\smallskip

\noindent\textit{Editorial scope.} Counted unit: the article's lemma lem:variational-principle (source0.tex lines 734--739). The article's complete original proof of it is delivered with it (source0.tex lines 740--755, one \texttt{proof} environment of 16 lines). No mathematics is written, repaired or re-derived: the statement and the proof are the article's own lines, in the article's own notation, and no retained line is substituted or re-flowed.  No further source-local context is needed: every cross-reference of every retained line resolves inside the two delivered spans.  Nothing was omitted from any retained span, so the package carries no omission record.  No retained line contains a citation, so the wrapper carries no bibliography and no bibliography entry.  No retained line contains a figure environment, an image-inclusion command or drawing code, so no image is delivered and the package declares no asset dependency.  Editorial formatting only, in addition: an AMS \texttt{amsart} preamble that restates the article's own declarations and macros used by the retained lines, namely the environments \texttt{lemma} on separate counters, together with the standard packages those lines need.  The retained environments print in this document as Lemma 1 in order of appearance, whereas the article prints the same items with the numbering of its own full document.  The complete native source of the article as distributed in the arXiv e-print for this exact version is bundled unchanged as \texttt{sources/194315/source0.tex}.
\begin{lemma}[Variational principle of the Fisher information] \label{lem:variational-principle} Let $\mc{C}^\infty_c(\R^n; \R^n)$ be the family of all compactly supported and smooth vector fields from $\R^n$ to $\R^n$, then for all $t \geq 0$,
\begin{align*}
    \jfunc(t) = \sup _{b \in \mc{C}_c^{\infty}(\R^n; \R^n)}\left\{- \frac{2}{n} \int p_t(y) \nabla \cdot b(y) \de y- \frac{1}{n}\int p_t(y) \norm{b(y)}^2 \de y\right\} .
\end{align*}
    
\end{lemma}

\begin{proof}
    By nonnegativity
    \begin{align*}
        \frac{1}{n}\int p(y)\norm{b(y) - \nabla \log p_t(y)}^2 \de y = \epsilon_b \geq 0,
    \end{align*}
    we immediately have $\jfunc(t) \geq \epsilon_b + \frac{2}{n} \int p_t(y) b(y) \cdot \nabla \log p_t(y) \de y - \frac{1}{n}\int p_t(y) \norm{b(y)}^2 \de y$. While $b$ is smooth and compactly supported, we can integrate by parts and obtain
    \begin{align*}
        \frac{2}{n} \int p_t(y) b(y) \cdot \nabla \log p_t(y) \de y = \frac{2}{n}  \int b(y) \cdot \nabla p_t(y) \de y = -  \frac{2}{n}  \int p_t(y) \nabla \cdot b(y) \de y.
    \end{align*}
    Here $\nabla \cdot b(y) := \sum_{i=1}^n \frac{\partial}{\partial y_i} b(y)$ denotes the divergence of $b(y)$. Thus taking supremum over all $b \in \mc{C}_c^{\infty}(\R^n; \R^n)$,
    \begin{align*}
        \jfunc(t) & \geq \sup _{b \in \mc{C}_c^{\infty}(\R^n; \R^n)}\left\{ \epsilon_b - \frac{2}{n} \int p_t(y) \nabla \cdot b(y) \de y- \frac{1}{n}\int p_t(y) \norm{b(y)}^2 \de y \right\} \\
        & \geq \sup _{b \in \mc{C}_c^{\infty}(\R^n; \R^n)}\left\{- \frac{2}{n} \int p_t(y) \nabla \cdot b(y) \de y- \frac{1}{n}\int p_t(y) \norm{b(y)}^2 \de y\right\}.
    \end{align*}
    The equality holds since $\epsilon_b$ can be arbitrarily small.
\end{proof}

\end{document}
